\section{MVP Design}
\label{sec:mvp}

This section defines the Minimum Viable Product for Timeline Compiler, identifying must-have features, phased delivery, risks, and complexity hotspots.

\subsection{MVP Philosophy}

The MVP must answer one question definitively: \emph{Can Timeline Compiler produce presentation-quality timelines from structured input, with deterministic output and theme-driven styling?}

Everything else — advanced features, edge cases, integrations — can follow once this core loop is validated. The MVP is not a prototype; it is a shippable product with limited scope.

\subsection{Feature Prioritization}

\subsubsection{Must Have (P0)}

These features are required for MVP launch:

\begin{description}
  \item[IR Parser and Validator] Parse YAML input conforming to the Timeline IR schema. Validate against JSON Schema. Report clear, actionable errors.
  
  \item[Core IR Elements] Support for:
  \begin{itemize}
    \item \texttt{timeline} root with \texttt{title} and \texttt{range}
    \item \texttt{tracks} (swimlanes)
    \item \texttt{activities} with labels, date ranges, status, and progress
    \item \texttt{milestones} with dates and labels
  \end{itemize}
  
  \item[Deterministic Layout Engine] Compute positions for all elements deterministically. Handle overlapping activities within tracks (stacking or compression).
  
  \item[SVG Output] Primary output format. Clean, well-structured SVG suitable for web embedding and further processing.
  
  \item[PDF Output] Print-quality vector output for documents and decks.
  
  \item[Basic Theme Support] At least one built-in theme (\texttt{default}). Theme defines colors, typography, spacing, and shapes.
  
  \item[CLI] Command-line interface with:
  \begin{itemize}
    \item \texttt{timeline render <input> -o <output>} — render to file
    \item \texttt{timeline validate <input>} — validate IR without rendering
    \item \texttt{--theme <name>} — select theme
    \item \texttt{--format <svg|pdf>} — output format (or infer from extension)
  \end{itemize}
  
  \item[Schema Documentation] Published JSON Schema for the IR, with human-readable documentation.
\end{description}

\subsubsection{Should Have (P1)}

These features significantly improve usability and should be included if feasible:

\begin{description}
  \item[PNG Output] Rasterized output for contexts requiring bitmap images.
  
  \item[Multiple Built-in Themes] At least 3 themes: \texttt{default}, \texttt{minimal}, \texttt{corporate}. Demonstrates theme system versatility.
  
  \item[Sections] Temporal divisions (e.g., quarters, phases) that span all tracks with visual demarcation.
  
  \item[Status Visualization] Color or icon mapping for status values (\texttt{planned}, \texttt{in-progress}, \texttt{complete}, \texttt{at-risk}).
  
  \item[Progress Bars] Visual representation of \texttt{progress} (0.0–1.0) within activity bars.
  
  \item[Today Marker] Vertical line at a specified ``today'' date.
  
  \item[npm Package] Core library published to npm for embedding.
  
  \item[Custom Theme Loading] Load themes from user-provided YAML/JSON files.
\end{description}

\subsubsection{Nice to Have (P2)}

These features are desirable but not blocking for MVP:

\begin{description}
  \item[PPTX Output] PowerPoint-compatible output. Complex due to PPTX format internals.
  
  \item[VS Code Extension] Live preview, syntax highlighting. High value but significant development effort.
  
  \item[Groups] Nested groupings within tracks for hierarchical organization.
  
  \item[Annotations] Callouts and emphasis markers for specific dates or activities.
  
  \item[Dependency Arrows] Visual arrows between activities (visual only, no scheduling).
  
  \item[Legend Generation] Auto-generated legends for status colors or track semantics.
  
  \item[Watch Mode] CLI watches input file and re-renders on change.
  
  \item[MCP Server] Agent integration. Important for long-term vision but not required for initial launch.
\end{description}

\subsubsection{Future (P3)}

These features are out of scope for MVP but part of the long-term vision:

\begin{description}
  \item[Theme Marketplace] Community-contributed themes.
  
  \item[Multi-file IR Composition] \texttt{\$ref} or include mechanisms for large timelines.
  
  \item[External Image Embedding] Logos, icons from external files.
  
  \item[Interactive SVG] Clickable regions, tooltips (breaks determinism constraints — careful design needed).
  
  \item[Hosted Rendering API] Cloud endpoint for rendering without local install.
  
  \item[Adapter Ecosystem] Tools that convert Jira, ADO, GitHub Projects to Timeline IR.
  
  \item[Localization] Date formats, RTL support, translated labels.
\end{description}

\subsection{MVP Scope Summary}

\begin{table}[h]
\centering
\caption{MVP Feature Summary}
\begin{tabularx}{\textwidth}{llX}
\toprule
\textbf{Priority} & \textbf{Label} & \textbf{Features} \\
\midrule
P0 & Must Have & IR parser, validation, tracks, activities, milestones, layout engine, SVG, PDF, CLI, 1 theme, schema docs \\
P1 & Should Have & PNG, 3 themes, sections, status colors, progress bars, today marker, npm package, custom themes \\
P2 & Nice to Have & PPTX, VS Code extension, groups, annotations, dependency arrows, legends, watch mode, MCP server \\
P3 & Future & Theme marketplace, multi-file IR, images, interactive SVG, hosted API, adapters, localization \\
\bottomrule
\end{tabularx}
\end{table}

\subsection{Phased Roadmap}

\subsubsection{Phase 1: Core Engine (MVP Launch)}

\textbf{Duration:} 6–8 weeks

\textbf{Deliverables:}
\begin{itemize}
  \item IR schema (JSON Schema) finalized and documented
  \item Parser and validator
  \item Layout engine (deterministic)
  \item SVG renderer
  \item PDF renderer (via SVG conversion or direct)
  \item CLI with \texttt{render} and \texttt{validate} commands
  \item Default theme
  \item Public repository, MIT license, README, contributing guide
\end{itemize}

\textbf{Exit Criteria:}
\begin{itemize}
  \item A 3-track, 10-activity roadmap renders to SVG and PDF
  \item Output is byte-identical across runs
  \item CLI works on macOS, Linux, Windows
\end{itemize}

\subsubsection{Phase 2: Polish and Usability}

\textbf{Duration:} 4–6 weeks (overlapping with Phase 1 completion)

\textbf{Deliverables:}
\begin{itemize}
  \item PNG output
  \item 2 additional themes (\texttt{minimal}, \texttt{corporate})
  \item Sections support
  \item Status and progress visualization
  \item Today marker
  \item npm package published
  \item Custom theme loading from file
\end{itemize}

\textbf{Exit Criteria:}
\begin{itemize}
  \item Users can author themes without modifying core code
  \item npm package installable and functional
\end{itemize}

\subsubsection{Phase 3: Ecosystem}

\textbf{Duration:} 8–12 weeks

\textbf{Deliverables:}
\begin{itemize}
  \item VS Code extension (syntax highlighting, live preview, export)
  \item MCP server (render tool for agents)
  \item PPTX output
  \item Watch mode in CLI
  \item Groups and annotations in IR
\end{itemize}

\textbf{Exit Criteria:}
\begin{itemize}
  \item AI agent can generate IR and invoke rendering
  \item VS Code extension in marketplace
\end{itemize}

\subsubsection{Phase 4: Scale and Community}

\textbf{Duration:} Ongoing

\textbf{Deliverables:}
\begin{itemize}
  \item Theme contribution workflow
  \item Adapter examples (Jira, GitHub, ADO)
  \item Documentation site
  \item Advanced IR features (dependency arrows, legends)
  \item Performance optimization for large timelines
\end{itemize}

\subsection{Risks and Mitigations}

\subsubsection{Technical Risks}

\begin{description}
  \item[PDF Generation Complexity] PDF output requires either converting SVG (quality/fidelity concerns) or direct PDF writing (complex). \emph{Mitigation:} Use proven libraries; accept minor fidelity differences in MVP; iterate.
  
  \item[Font Determinism] Different systems have different fonts. Varying fonts break determinism. \emph{Mitigation:} Require explicit font specification in themes; bundle fallback fonts; document font requirements.
  
  \item[Layout Edge Cases] Overlapping activities, long labels, dense timelines may produce poor layouts. \emph{Mitigation:} Start with simple layout rules; document limitations; accept that MVP handles ``typical'' cases, not all.
  
  \item[Cross-Platform Rendering Differences] Floating-point differences across platforms may affect layout. \emph{Mitigation:} Use integer math or fixed-point where possible; test on multiple platforms early.
\end{description}

\subsubsection{Adoption Risks}

\begin{description}
  \item[Learning Curve] Users must learn the IR format. \emph{Mitigation:} Excellent examples, quick-start guide, VS Code extension with auto-complete.
  
  \item[Theme Quality] If built-in themes are unappealing, users won't adopt. \emph{Mitigation:} Invest in theme design; consult with visual designers; study exemplary slides.
  
  \item[Missing Feature X] Users may expect Gantt-like features (dependencies, resources) that are out of scope. \emph{Mitigation:} Clear documentation of scope boundaries; position as ``communication tool, not PM tool.''
\end{description}

\subsubsection{Complexity Hotspots}

\begin{description}
  \item[PPTX Output] The PPTX format (Office Open XML) is complex. Producing valid, editable PPTX files is non-trivial. \emph{Mitigation:} Defer to Phase 3; use established libraries; accept limited editability initially.
  
  \item[Theme System] Balancing flexibility with simplicity is hard. Too simple: users can't customize. Too complex: themes become a programming language. \emph{Mitigation:} Start minimal; extend based on feedback; keep theme schema documented.
  
  \item[Layout Engine] Deterministic layout that handles edge cases (long text, overlap, many tracks) is complex. \emph{Mitigation:} Define clear layout rules; document limitations; iterate.
\end{description}

\subsection{Smallest Viable Version}

If resources are constrained, the absolute smallest viable version is:

\begin{itemize}
  \item IR schema for \texttt{tracks}, \texttt{activities}, \texttt{milestones}
  \item YAML parser and validator
  \item SVG output only
  \item CLI with \texttt{render} command only
  \item Single hardcoded theme
  \item No npm package (CLI only)
\end{itemize}

This version proves the core thesis: structured input to presentation-quality output, deterministically. Everything else is enhancement.

\subsection{Success Metrics}

\begin{description}
  \item[Adoption] GitHub stars, npm downloads, VS Code extension installs.
  
  \item[Quality] Issue count, bug severity, user-reported rendering defects.
  
  \item[Determinism] CI tests that verify byte-identical output across runs and platforms.
  
  \item[Agent Usage] Number of agents using MCP server (post-Phase 3).
  
  \item[Community Contribution] Theme contributions, adapter contributions, PRs.
\end{description}
